\documentclass[11pt]{amsart}
\usepackage[T1]{fontenc}
\usepackage{lmodern,amsmath,amssymb,mathtools,microtype,booktabs}
\usepackage[margin=1.15in]{geometry}
\usepackage[hidelinks]{hyperref}
\newtheorem{theorem}{Theorem}[section]
\newtheorem{proposition}[theorem]{Proposition}
\newtheorem{lemma}[theorem]{Lemma}
\newtheorem{corollary}[theorem]{Corollary}
\theoremstyle{definition}\newtheorem{definition}[theorem]{Definition}
\theoremstyle{remark}\newtheorem{remark}[theorem]{Remark}
\newcommand{\R}{\mathbb R}
\newcommand{\Z}{\mathbb Z}
\newcommand{\E}{\mathbb E}
\newcommand{\Prob}{\mathbb P}
\newcommand{\one}{\mathbf 1}
\newcommand{\dd}{\,\mathrm d}
\DeclareMathOperator{\dist}{dist}
\DeclareMathOperator{\covol}{covol}
\DeclareMathOperator{\area}{area}
\DeclareMathOperator{\diam}{diam}
\numberwithin{equation}{section}
\title[Scalar collision laws]{Scalar collision laws and recognition of periodic dispersing billiards}
\author{Qian Qi}
\date{October 4, 2026}
\subjclass[2020]{37D50, 52A10, 60J10, 62G05, 68Q32}
\keywords{Dispersing billiards, reciprocal collision commands, active boundary estimation,
bit complexity, periodic recognition, conditional reconstruction}
\hypersetup{pdftitle={Scalar collision laws and recognition of periodic dispersing billiards},pdfauthor={Qian Qi}}
\begin{document}
\begin{abstract}
We reconstruct periodic dispersing tables from localized reciprocal
collision commands, each returning one bit. Solid starts and free misses
both return zero and remain in the denominator. A fixed-depth local
Bellman computation converts pooled reciprocal means into approximate
membership without a supplied free point or component boundary. After
a fixed coarse acquisition, adaptive radial searches and polynomial
reconstruction recover boundaries with $C^2$ error $\nu$. On bounded
$C^{6,\beta}$ classes with known geometric and nonperiod-patch margins,
the number of attempted preparations is at most
$C\nu^{-1/(4+\beta)}\log(C/\nu)\log(C/(\nu\delta))$
at confidence $1-\delta$. The localization and coupled calibration scale
is $\nu^{(6+\beta)/(4+\beta)}$. A physical packing and a binary-transcript
bound require $c\nu^{-1/(4+\beta)}$ bits on a fixed nondegenerate subclass,
even with its period and pose given. Thus the power of the attempted-bit
complexity is determined, with a logarithmic gap. Protected-patch tests
also recover the primitive orbit count and exact rational period
relations; Euclidean generators remain estimates. The result is an
active, calibrated input-output theorem, not a passive count or spectral
rigidity theorem. The exact reciprocal inverse and all preceding results
are retained in the accompanying source package.
\end{abstract}
\maketitle
\section{The observation and the main results}
\label{sec:setting}
The number of outcomes of a sensor and the information in its controls
are different resources. Here every preparation returns one bit, but
the controller chooses its spatial input. Our question is how many such
preparations are needed for geometric reconstruction, once their input
precision is specified. The reciprocal identity below reduces collision
observations to an occupation difference. The main point is that this
reduction can be evaluated at isolated, adaptively chosen locations;
a fine grid throughout the observation aperture is unnecessary.

\subsection{Physical class and loss}
Fix $0<\beta\le1$ and put $s=6+\beta$. A table is the nonempty union
\begin{equation}\label{eq:presentation}
 \mathcal O=\bigcup_{i=1}^r(C_i+L\Z^2),\qquad 1\le r\le r_0.
\end{equation}
This presentation has exactly $r$ component orbits. It need not be
primitive. The constants in the following bounds are fixed and known:
\begin{equation}\label{eq:priors}
 \|L\|,\|L^{-1}\|\le L_0,\quad |\det L|\ge V_0>0,\quad
 \diam C_i\le D_0,\quad \dist(C,D)\ge d_0>0
\end{equation}
for distinct physical components. All components are strictly convex,
with curvature in $[\kappa_-,\kappa_+]\subset(0,\infty)$ and centered
support functions bounded by $M$ in $C^{6,\beta}(\R/2\pi\Z)$.
Reflections, repeated shapes and nonprimitive presentations are allowed.
Write $\mathcal B$ for this class, $\Pi=\{w:\mathcal O+w=\mathcal O\}$
for its full period group, $z_C$ for a component's Steiner point, and
$p_C$ for its centered support function. All norms on finite-dimensional
spaces are fixed once and for all.

A known positive nonperiod-patch margin $\eta$ will be required only
for uniform finite discrete decisions. Its precise definition is
\eqref{eq:patch-margin}. We denote the resulting subclass by
$\mathcal B_\eta$. The geometric loss is the maximum $C^2$ error of
matched component support functions in a protected laboratory patch,
together with the error of a matched primitive period basis and free
area. Representatives are chosen in a common protected patch. The
basis is matched to a true basis, not required to be a continuously
chosen reduced basis. Exact rational relations and the primitive orbit
count are stated separately from this real-valued loss.

\subsection{Commands and their accounting}
For a nominal displacement $a$, let $B_a(q)$ equal one if a free start
at $q$ has a first collision on the unreflected segment $[q,q+a]$.
It equals zero for a solid start and for a free miss. Both are attempted
preparations and enter the denominator. Boundary conventions on
zero-volume sets are immaterial. Fix an even smooth probability kernel
$\kappa$ supported in the unit disk, and set
\[
 f_{x,\sigma}(q)=\sigma^{-2}\kappa((q-x)/\sigma),\qquad
 u_\sigma=\kappa_\sigma*\one_{\mathcal O}.
\]
Choose $\sigma_0,t>0$ with $t+2\sigma_0<d_0$, and prescribe the compass
law
\begin{equation}\label{eq:compass}
 \rho=\tfrac14(\delta_{te_1}+\delta_{-te_1}
                         +\delta_{te_2}+\delta_{-te_2}).
\end{equation}
A forward command draws $q\sim f_{x,\sigma}$ and $a\sim\rho$
independently and attempts $(q,a)$. A reverse command draws that same
nominal joint law and attempts $(q+a,-a)$. Separate fresh preparations
are used; pairing the same physical particle is unnecessary. Only the
two pooled means are used, not means resolved by direction. The joint
reverse law, rather than only its marginal direction law, is prescribed.

Each actual preparation may be coupled to its nominal one with position,
duration and angle errors bounded by $\ell,\tau,\alpha$, respectively.
It remains a unit-speed straight flight until its first impact. The
bounds hold conditionally on past commands, and preparations at a given
command are fresh and conditionally independent. Section~\ref{sec:queries}
proves the bias bound $C(\ell+\tau+t\alpha)/\sigma$. Calibration is a
resource of the controller, not an estimate inferred from the bits.
The launch aperture is fixed by the prior and does not grow as accuracy
increases. It includes the translated reverse starts and the kernel
supports. No free point, obstacle label, impact location, impact angle,
collision time, primitive cell or area normalization is supplied.

\begin{theorem}[Adaptive scalar reconstruction]\label{thm:adaptive}
Fix the numerical bounds of $\mathcal B_\eta$ and the kernel above.
There are $C,c,\nu_0>0$ such that, for $0<\nu<\nu_0$ and
$0<\delta<1/4$, a finite adaptive experiment with the reciprocal
compass commands reconstructs a smooth strictly convex periodic
presentation with geometric loss at most $C\nu$, with probability
at least $1-\delta$. It recovers the true primitive orbit count and
exact rational relations of a true period-generating list in a true
independent-pair basis. The Euclidean coordinates of these vectors
are estimates.

After a prior-dependent coarse acquisition, the finest localization
scale and sufficient coupled calibration are
\begin{equation}\label{eq:precision}
 \sigma=c\nu^{s/(s-2)},\qquad \ell+\tau+t\alpha\le c\sigma.
\end{equation}
The number of nominal command centers, counted with repetitions if
necessary, and the number of attempted preparations satisfy
\begin{align}
 J_\nu&\le C\nu^{-1/(s-2)}\log(C/\nu),\label{eq:centers}\\
 N_\nu&\le C\nu^{-1/(s-2)}\log(C/\nu)
                               \log(C/(\nu\delta)).\label{eq:attempts}
\end{align}
Every solid start and free miss is charged in $N_\nu$. All commands lie
in one fixed prior-controlled aperture. The reconstruction uses finite
local computations, radial bisections and finite polynomial data, not
a search over an infinite-dimensional class of physical candidates.
Its constants include the exit-time and patch-margin bounds.
\end{theorem}

Here and below a finite smooth partition of unity, with its derivatives,
is part of the fixed numerical representation. Certified evaluations
to the reserved tolerances suffice. The theorem counts preparations
and centers, not apparatus travel or bit operations for elementary
functions. Known $s$ is used; adaptivity refers to the locations of
commands, not adaptation to unknown smoothness.

\begin{theorem}[A necessary number of binary observations]
\label{thm:lower}
For each fixed $0<\beta\le1$ there is a fixed choice of the bounds
above and a nondegenerate subclass $\mathcal P\subset\mathcal B_\eta$
with the following property. Any adaptive randomized procedure whose
only table-dependent outputs are at most $N$ bits, and which recovers
the component boundaries with $C^2$ error at most $\nu$ with probability
at least $3/4$ uniformly on $\mathcal P$, must satisfy
\begin{equation}\label{eq:lower}
                         N\ge c\nu^{-1/(s-2)}.
\end{equation}
This holds even when the lattice, its primitive orbit count, component
correspondences and laboratory pose are given, and implementation is
exact. Input choices may have arbitrary precision and depend on all
previous bits. Independent controller randomness is allowed.
\end{theorem}

The construction of $\mathcal P$ is given in Section~\ref{sec:lower}.
The assertion is not a lower bound for every subclass, in particular
not for a singleton or a finite parametric family. On any bounded
class containing this construction, \eqref{eq:attempts} and
\eqref{eq:lower} identify the power of the attempted-bit complexity.
A logarithmic gap remains. Neither result prices input precision or
converts an active experiment into a passive billiard invariant.

\section{From reciprocal collisions to local occupation queries}
\label{sec:queries}
This section supplies the observation reduction used by every adaptive
query. The random walk is computational: its exit time is not an
observed collision record or an apparatus stopping signal.

\begin{lemma}[Reciprocal balance and calibration]\label{lem:balance}
With $\chi=\one_{\mathcal O}$, outside null boundary sets,
\begin{equation}\label{eq:balance}
 B_a(q)-B_{-a}(q+a)=\chi(q+a)-\chi(q).
\end{equation}
Consequently the difference $g_\sigma$ of the two pooled nominal means
satisfies
\begin{equation}\label{eq:forcing}
 g_\sigma=(T-I)u_\sigma,\qquad
 Tv(x)=\tfrac14\sum_{a\in\{\pm te_1,\pm te_2\}}v(x+a).
\end{equation}
For either command, coupled errors of size
$r=\ell+\tau+t\alpha\le\sigma$ change its mean by at most $Cr/\sigma$.
The constant is uniform over the protected aperture, components and
nominal directions, including tangencies.
\end{lemma}
\begin{proof}
If both endpoints are free, a segment and its reverse either both hit
or both miss. If exactly one endpoint is solid, the direction starting
there gives zero and the other direction hits. If both are solid both
bits are zero. This proves \eqref{eq:balance}. Averaging in $q,a$ and
commuting translation with convolution proves \eqref{eq:forcing}.

For a convex body put $S_a(C)=C+[-a,0]$. The hit bit is
\[
 B_a(q)=\one_{\cup_C S_a(C)}(q)-\chi(q).
\]
The displacement error is at most $\tau+t\alpha$. Support inequalities
show that membership can change only if the nominal start lies in
an $r$-tube of $\partial C$ or $\partial S_a(C)$, including the position
error. The part of a convex boundary tube in a disk of radius $\sigma$
has area at most $C(\sigma+r)r$. One proof integrates boundaries of
outer parallel bodies and inner erosions by the coarea formula; their
portions in the disk have length at most $2\pi\sigma$, by monotonicity
of the perimeter of convex sets. Degenerate levels follow by a limit.
Only a bounded number of components contribute: each has a point in
a disk of radius $\sigma+t+r$, and such points from distinct components
are $d_0$-separated. Multiplication by the density bound
$\sigma^{-2}\|\kappa\|_\infty$ proves the assertion. The tube containment
is pointwise for every allowed error, so the coupling need not have
mean zero. An exceptional coupling event adds its probability.
\end{proof}

Put $D=D_0+2\sigma_0$, $H=(D+t)^2/t^2$ and
$L_* =\max(1,\lceil2H\rceil)$. The positive set of $u_\sigma$ lies in
the $\sigma$-enlarged components. They have diameter at most $D$ and
mutual distance greater than $t$.

\begin{lemma}[Stopped inverse in one aperture]\label{lem:stopped}
For the compass walk $X_n=x+a_1+\cdots+a_n$ and
$\tau_u=\inf\{n\ge0:u_\sigma(X_n)=0\}$ one has
\begin{equation}\label{eq:survival}
 \sup_x\E_x\tau_u\le H,\qquad
 \sup_x\Prob_x(\tau_u>n)\le2^{-\lfloor n/L_*\rfloor}.
\end{equation}
Let a bounded target region $U$ contain the complete protected bodies
and their search brackets. Choose a fixed bounded $W$ containing
$U+\overline B_{D+t+1}$. On $W$ use the killed transition
$T^Wv=T(\one_Wv)$ and set iterates to zero outside $W$. For
$V_0=0$, $V_{n+1}=\max(0,T^WV_n-g_\sigma)$,
\begin{equation}\label{eq:bellman-error}
 0\le u_\sigma-V_n\le2^{-\lfloor n/L_*\rfloor}\quad\hbox{on }U.
\end{equation}
Uniform forcing error $\xi$ and numerical update error $\zeta$ add
at most $n(\xi+\zeta)$ if approximate iterates are clipped to $[0,1]$.
\end{lemma}
\begin{proof}
Until its first zero, a path stays in one enlarged component, since
one step cannot reach another. For $S=\tau_u\wedge n$,
$|X_S-x|\le D+t$. The bounded stopping identity for the martingale
$|X_n-x|^2-nt^2$ gives $t^2\E S\le(D+t)^2$. Monotone convergence
proves the mean bound. Markov's inequality at $L_*$ and the Markov
property at successive blocks prove the geometric tail.

On $W$, $T^Wu_\sigma\le Tu_\sigma$, so $0\le V_n\le u_\sigma$
and the iterates vanish on its zero set. Starting in $U$, the complete
containing component and every first-exit location lie in $W$. Along
this stopped path the killed transition has no boundary discrepancy.
For $e_n=u_\sigma-V_n$, the recursion gives $e_{n+1}\le Te_n$ before
exit, and hence
\[
 e_n(x)\le\E_x[u_\sigma(X_n)\one_{\{\tau_u>n\}}].
\]
This proves \eqref{eq:bellman-error}. Maximum, averaging and clipping
are nonexpansive in supremum norm; induction proves the error bound.
No periodic wrap or renormalization of a boundary row is used.
\end{proof}

The stopped argument also compares two physical occupations with
unrelated zero sets. Indeed, for $w=u-\widetilde u$ and
$h=(T-I)w$, bounded stopping at $\tau_u\wedge n$ gives
\[
 w(x)=\E_xw(X_{\tau_u\wedge n})
                 -\E_x\sum_{j<\tau_u\wedge n}h(X_j).
\]
At exit the terminal value is $-\widetilde u\le0$. On survival it
is at most one. Letting $n$ increase and then interchanging $u$ and
$\widetilde u$ yields $\|u-\widetilde u\|_{U,\infty}
\le H\|(T-I)(u-\widetilde u)\|_{W,\infty}$. Thus exact reciprocal
functionals at all sufficiently small scales determine the local
occupation and, by approximate-identity convergence and regular
closedness, the physical set. This retained exact conclusion requires
neither analyticity nor a supplied zero set. The martingale and killed
Green mechanisms are classical; see \cite{LawlerLimic}.

\begin{proposition}[A finite local query from pooled bits]
\label{prop:query}
For any adaptively chosen $x\in U$ and $\sigma\le\sigma_0$, an
occupation estimate with error less than $1/8$ can be computed using
only a prior-bounded number of reciprocal command centers in $W$.
At conditional confidence $1-\varepsilon$ it uses at most
$C\log(C/\varepsilon)$ attempted preparations. A sufficient calibration
is $\ell+\tau+t\alpha\le c\sigma$. Thresholding the estimate at
$1/2$ gives the correct membership label outside the $\sigma$-tube
of the component boundaries. Its label inside that tube is unrestricted.
\end{proposition}
\begin{proof}
Fix $n_*=6L_*$. To evaluate $V_{n_*}(x)$, its dependency locations
are exactly among
\begin{equation}\label{eq:diamond}
 x+t(i,j),\qquad |i|+|j|\le n_*.
\end{equation}
Intersect these with $W$ and put zero at locations outside it. Evaluate
the recursion backwards through shrinking diamonds, initializing
$V_0=0$. Induction on the layer shows equality with the full killed
iterate at the target. There are $O(n_*^2)$ centers and $O(n_*^3)$
updates, both independent of $\sigma$. No fine background spatial grid,
interpolation of the transition, or growing aperture is involved.

At each center estimate each pooled mean to error at most
$1/(1024n_*)$, and impose the same bound on its calibration bias.
The forcing error is then at most $1/(256n_*)$. Reserve update error
at most $1/(256n_*)$. Lemma~\ref{lem:stopped} gives total error at most
$1/64+1/256+1/256<1/8$. Hoeffding's inequality and a union bound over
the fixed finite number of means give the preparation bound. The
constant includes $n_*$ and may be large. The proof holds conditional
on past data because the new attempts are fresh. Inside a component
at boundary distance greater than $\sigma$ the average is one;
outside all enlarged components it is zero. Thresholding therefore
has the stated property. A membership label is a computed conclusion
from these bits, not an extra response furnished by the sensor.
\end{proof}

The full launch region may be taken to contain
$W+\overline B_{t+\sigma_0+1}$, which includes reverse positions and
kernel supports. For at most $Q$ adaptive queries, allocating conditional
failure probability $\delta/Q$ to each gives joint correctness with
probability $1-\delta$. It suffices to know a deterministic upper bound
on $Q$; repeated centers may be resampled. No independence of the
adaptively selected locations themselves is asserted.

\section{Adaptive recovery of complete component boundaries}
\label{sec:boundary}
We first acquire a coarse, fixed-resolution description. Its purpose
is to find interior centers and isolated radial brackets, not to supply
the final accuracy. This avoids both a supplied free reference point
and a bisection that might accidentally encounter a different obstacle.

\begin{lemma}[Coarse hulls and interior centers]\label{lem:coarse}
A prior-dependent finite number of queries from
Proposition~\ref{prop:query}, at one fixed scale $\sigma_c>0$, recovers
all complete components needed in the protected patch to Hausdorff
error $e_c\le3\sigma_c$. It supplies, for each such component $C$, a
polygon $H$ and a computable center $o$ for which
\begin{equation}\label{eq:inner-ball}
 \overline B_r(o)\subset C\cap H\subset\overline B_{R_1}(o)
\end{equation}
with fixed positive $r,R_1$. For every direction $v$, it supplies a
radial interval of fixed length at most $C e_c$ containing the boundary
point of $C$ on $o+\R_+v$. The interval is disjoint from the
$\sigma$-enlargements of all other components for every sufficiently
small fine scale $\sigma$.
\end{lemma}
\begin{proof}
The curvature upper bound supplies an interior rolling radius
$r_* =1/\kappa_+$. In support coordinates this follows from
$p+p''\ge r_*$: subtracting the support function of the radius-$r_*$
disk leaves a support function with nonnegative curvature measure.
In particular $C=C_{-2\sigma_c}+2\sigma_c\overline B$ when
$2\sigma_c<r_*$. Choose $\sigma_c<\min(r_*/8,d_0/40)$ and a square
grid of covering radius $\rho_c\le\sigma_c/8$ on a surrounding
square with collar larger than $2D_0+2$.

At its nodes retain the estimated occupations at least $1/2$, join
retained nodes at distance at most $4\sigma_c$, and take a convex hull
for each connected cluster. Every retained node is within $\sigma_c$
of a unique body; all nodes of $C_{-\sigma_c}$ are retained. For any
retained node assigned to $C$, choose a point in $C_{-2\sigma_c}$ at
distance at most $3\sigma_c$, and then a nearest grid node. That node
is retained and joined to the first. Convexity connects any two such
deep nodes by a segment; partitioning it into lengths at most $\rho_c$
and replacing partition points by nearest nodes yields a retained
path with successive lengths at most $3\rho_c$. Different bodies
cannot join because $d_0-2\sigma_c>4\sigma_c$. The hull thus lies in
$C+\sigma_c\overline B$, whereas every point of $C$ is within
$2\sigma_c+\rho_c$ of a retained deep node. This proves the error.
Discard clusters whose hulls are within $D_0+1$ of the outer square
boundary. Fragments cannot survive this cutoff, and the chosen collar
protects every required complete component.

Compute a largest inscribed disk of $H$ by its finite half-plane
inequalities, to any fixed reserved precision. Since
$d_H(H,C)\le e_c$, the support inequalities imply
$C_{-e_c}\subset H$ and $H_{-e_c}\subset C$. A disk of radius $r_*$
in $C$ therefore gives an inscribed disk of radius at least $r_*-e_c$
in $H$. The computed center, with the radii reduced by the error and
numerical reserve, satisfies \eqref{eq:inner-ball} for a fixed $r>0$.
Diameter bounds supply $R_1$.

Write $R_H(v),R_C(v)$ for the radial functions about $o$. Since both
bodies contain $\overline B_r(o)$,
\[
 H+e_c\overline B\subset o+(1+e_c/r)(H-o),
\]
and the analogous inclusion holds for $C$. The two Hausdorff
inclusions give $|R_H(v)-R_C(v)|\le C e_c$ uniformly. The polygonal
radial value is computable from its sides. Use a slightly enlarged
bracket about it, allowing numerical error. By decreasing the fixed
$\sigma_c$, its length and every point's distance to $C$ are less
than $d_0/4$. Separation excludes every other component from the
bracket and its sufficiently small kernel neighborhood.
\end{proof}

\begin{lemma}[Bisection with an indeterminate boundary layer]
\label{lem:bisection}
On each bracket in Lemma~\ref{lem:coarse}, the derived label of
Proposition~\ref{prop:query} may be arbitrary, and may depend on the
past, only when
\begin{equation}\label{eq:radial-ambiguity}
                       |a-R_C(v)|\le A\sigma
\end{equation}
for a fixed $A$. After $k$ midpoint tests, ordinary bisection returns
an estimate with error at most $A\sigma+2^{-k-1}w_0$, where $w_0$
is the initial bracket length. In particular $O(\log(C/\sigma))$
queries recover $R_C(v)$ to $O(\sigma)$.
\end{lemma}
\begin{proof}
Translate $o$ to zero. The inner ball implies
\[
 (1-\sigma/r)C\subset C_{-\sigma},\qquad
 C+\sigma\overline B\subset(1+\sigma/r)C.
\]
Together with $R_C\le R_1$, these show that outside
\eqref{eq:radial-ambiguity} the label is the true radial inside/outside
label. Other components were excluded from the bracket. A label one
at radius $a$ implies $a\le R_C+A\sigma$; a label zero implies
$a\ge R_C-A\sigma$. If the current endpoints are $l,u$, midpoint
updates consequently preserve the relaxed invariant
$R_C\in[l-A\sigma,u+A\sigma]$. This does not require $R_C$ to remain
in $[l,u]$, or the noisy labels to be monotone. The width halves at
every step. Its midpoint has the asserted error.
\end{proof}

\begin{lemma}[Radial regularity and stable interpolation]
\label{lem:radial}
Under \eqref{eq:inner-ball} and the physical support priors, the
$2\pi$-periodic radial function $R_C(\varphi)$ has a uniformly bounded
$C^{6,\beta}$ norm. Let $h=2\pi/m$ for an integer $m\ge9$. From
radial values at the $m$ equally spaced angles, each with error at most
$e$, one can construct a smooth function $\widehat R$ such that
\begin{equation}\label{eq:radial-interpolation}
 \|\widehat R-R_C\|_{C^j}
           \le C(h^{s-j}+e h^{-j}),\qquad 0\le j\le3.
\end{equation}
For $e\le C_0h^s$ and sufficiently small $h$, the curve
$o+\widehat R(\varphi)(\cos\varphi,\sin\varphi)$ is the boundary of
a smooth strictly convex body $\widehat C$, and
\begin{equation}\label{eq:geometric-interpolation}
 d_H(\widehat C,C)\le Ch^s,\qquad
 \|p_{\widehat C}^{\rm lab}-p_C^{\rm lab}\|_{C^2}\le Ch^{s-2}.
\end{equation}
Here uncentered support functions are compared in the same laboratory
frame; centering preserves these orders.
\end{lemma}
\begin{proof}
Use the support function $p$ relative to $o$, so $p\ge r$. In normal
angle $\theta$, the boundary is $p n+p' n^\perp$, its radial angle is
\[
 \varphi=\theta+\arctan(p'/p),\qquad
 \frac{\dd\varphi}{\dd\theta}
             =\frac{p(p+p'')}{p^2+(p')^2}>c>0.
\]
The inverse $\theta(\varphi)$ is uniformly $C^{5,\beta}$. Although
the expression $R=\sqrt{p^2+(p')^2}$ initially gives only five
ordinary derivatives, differentiation in radial angle cancels the
curvature factor:
\begin{equation}\label{eq:radial-cancellation}
                   R'=R(p'/p)\circ\theta.
\end{equation}
Its right side is uniformly $C^{5,\beta}$. Thus $R$ is uniformly
$C^{6,\beta}$, without losing a derivative from the support prior.

At each angular node interpolate the seven consecutive sampled
values by a polynomial of degree at most six in a local lift of the
angle. Combine these polynomials with a fixed translated smooth
partition of unity supported on neighborhoods of size $O(h)$.
On each such neighborhood the normalized Vandermonde inverse is
bounded. A degree-six Taylor polynomial of the true function is
reproduced exactly, and its remainder and derivatives through order
three are $O(h^{s-j})$. Errors in the values contribute $O(eh^{-j})$.
Derivatives falling on the partition have size $O(h^{-j})$ and are
absorbed by the corresponding lower-order remainders. The overlap
number is fixed. This proves \eqref{eq:radial-interpolation}, including
across the angular cut by periodic lifts.

A positive radial function is strictly convex when
\[
       F=R^2+2(R')^2-RR''>0.
\]
For the true curve this expression has a uniform positive margin,
by the inradius, curvature and diameter bounds. It persists for
$\widehat R$ by the $C^2$ estimate. The two radial graphs are uniformly
$O(h^s)$-close, which proves the Hausdorff bound.

To check the stronger support norm, set
$\theta=\varphi-\arctan(R'/R)$, now viewing the normal angle as a
function of radial angle. Its derivative is
$F/(R^2+(R')^2)>c$. At corresponding angles the support value and
radius of curvature are
\[
 p=\frac{R^2}{\sqrt{R^2+(R')^2}},\qquad
 p+p''=\frac{(R^2+(R')^2)^{3/2}}{F}.
\]
These expressions, and the tangential support derivative, depend
Lipschitzly on $R,R',R''$ on the bounded set with these margins.
The inverse normal-angle maps differ by $O(\|\widehat R-R\|_{C^1})$.
Their composition changes the true expressions by the same order,
using the uniform $C^3$ bound; the reconstructed functions also have
such a bound by \eqref{eq:radial-interpolation}. Thus the support
error through two derivatives is $O(h^{s-2})$. Translation to the
laboratory frame adds the same first harmonic to both supports. The
Steiner point is a bounded linear functional of the support, so its
error is $O(h^s)$ and centering preserves the conclusions.
\end{proof}

\begin{proposition}[Adaptive patch reconstruction]\label{prop:patch}
For a fixed protected patch, complete component boundaries can be
recovered with the errors in \eqref{eq:geometric-interpolation} using
at most $C h^{-1}\log(C/h)$ local occupation queries. The finest
localization is $\sigma=c h^s$. With probability at least $1-\delta$
the attempted-bit cost is at most
\begin{equation}\label{eq:patch-cost}
           C h^{-1}\log(C/h)\log(C/(h\delta)).
\end{equation}
The number of components, the initial coarse queries and every
local Bellman stencil size are bounded solely by the fixed prior.
\end{proposition}
\begin{proof}
On the simultaneous accuracy event, obtain the hulls and centers in
Lemma~\ref{lem:coarse}. Separation and the fixed aperture bound their
number. For each body use $m$ equally spaced radial directions and
their isolated coarse brackets. Lemma~\ref{lem:bisection} uses
$O(\log(C/h))$ midpoint queries per direction to return values with
error $e\le C\sigma$. Lemma~\ref{lem:radial} gives the claimed
geometric estimates. The total number of queries is bounded by a
deterministic $Q_h=C(1+h^{-1}\log(C/h))$. Allocate failure probability
$\delta/Q_h$ to every query, including the coarse ones, before the
experiment. Conditional confidence and the union bound in
Proposition~\ref{prop:query} give the simultaneous event. Its
preparation cost is \eqref{eq:patch-cost}. Locations chosen using
coarse or previous fine data do not alter this argument.

For a finite numerical realization, polygon intersections, midpoint
locations and radial values are enclosed with error $O(\sigma)$.
The estimate \eqref{eq:radial-interpolation} absorbs those value errors.
Represent the output by its finite polynomial coefficients and the
fixed smooth partition functions. They, their derivatives and the
normal-angle maps can be evaluated to any positive prescribed error
by finite operations on compact intervals with the stated derivative
bounds. This constructs an actual smooth convex output, without
assigning a false smoothness bound to the initial polygons. No
infinite-dimensional candidate minimization is performed.
\end{proof}

\section{Period recognition from the reconstructed patch}
\label{sec:periods}
This section retains the protected-patch argument for repeated shapes.
It acts on the acquired geometry, not on an assumed list of integer
copy labels. We include the proof to specify exactly which discrete
information the adaptive experiment recovers.

Put $B=1+L_0$ and $Q_a=[-a,a]^2$. The aperture used above is chosen
large enough to recover complete components with centers in $Q_{10B}$,
together with their protective collars. For two components define
\[
 d_*(C,D)=\|z_C-z_D\|_\infty+\|p_C-p_D\|_\infty,
\]
where the centered support functions are compared in the laboratory
directions, without an independent rotation. Define the patch defect
\begin{equation}\label{eq:defect}
 \mathcal D(w)=\max_{C:z_C\in Q_B}\ \max_{\epsilon\in\{-1,1\}}
                     \inf_D d_*(C+\epsilon w,D).
\end{equation}
The infimum ranges over physical components. Distant components cannot
minimize it. The known nonperiod margin is the requirement
\begin{equation}\label{eq:patch-margin}
 \begin{gathered}
 w=z_D-z_C,\quad z_C\in Q_{3B},\ z_D\in Q_{5B},\quad w\notin\Pi
       \quad\Longrightarrow\quad\mathcal D(w)\ge\eta>0.
 \end{gathered}
\end{equation}
This is a margin for whole-patch translations, not a separation between
all individual shapes. Every fixed table in $\mathcal B$ has some
positive such margin, since there are finitely many candidate pairs.
The margin is not claimed to be uniformly inferable from finite bits.

\begin{lemma}[Exact patch recognition]\label{lem:recognition}
One has $\mathcal D(w)=0$ if and only if $w\in\Pi$. For
$\|w\|_\infty\le8B$, this zero test uses only components centered in
$Q_B$ and $Q_{9B}$. The group $\Pi$ is a lattice with
\begin{equation}\label{eq:index-bound}
 [\Pi:L\Z^2]\le r_0,\qquad\covol\Pi\ge V_*=V_0/r_0.
\end{equation}
Choose any component centered in $Q_{2B}$ as a root. Its center
differences to components centered in $Q_{4B}$ that pass the zero
test generate $\Pi$ over $\Z$.
\end{lemma}
\begin{proof}
Each orbit modulo $L\Z^2$ has a representative centered in
$L[-1/2,1/2]^2\subset Q_B$. A zero defect gives both translates of
every such representative as physical components. Translating by
$L\Z^2$ gives $\mathcal O+w\subset\mathcal O$ and
$\mathcal O-w\subset\mathcal O$; together these imply equality.
The converse is immediate. The centers of a matching pair in the
bounded test are in $Q_{9B}$.

Map a coset of $\Pi/L\Z^2$ to the component orbit of the translated
root. This is injective: a nonempty compact body is not invariant
under a nonzero translation. Thus the quotient has at most $r$ elements.
A finite extension of a full-rank lattice is a lattice, giving
\eqref{eq:index-bound}. The two columns of $L$, and representatives
of all these period cosets, can be chosen with norm less than $B$
in the maximum norm. Translating the root by these vectors gives
targets in $Q_{3B}$, strictly within the cutoff $Q_{4B}$. They are
included in the accepted list and generate $\Pi$. Conversely each
accepted vector is a true period.
\end{proof}

\begin{lemma}[Stable discrete decisions]\label{lem:locking}
Suppose the protected complete components have matched Hausdorff
error at most $e$. On $\mathcal B_\eta$, for sufficiently small $e$
depending on the full prior, finite tests recover the true primitive
orbit count and exact rational relations of a true accepted generating
list in a true independent-pair basis. They recover a matched period
basis and primitive free area to error $O(e)$.
\end{lemma}
\begin{proof}
Hausdorff error $e$ gives Steiner-center error at most $2e$ and
centered-support error at most $3e$. Use approximate source centers
in $Q_{2B}$, a root there and root targets in $Q_{4B}$. For a measured
difference $\widehat w$, its corresponding true difference has error
at most $4e$. Test both signs on all selected sources, minimizing over
targets with approximate centers in $Q_{10B}$. The center contribution
to each comparison changes by at most $8e$ and the support contribution
by at most $6e$. True periods therefore have measured defect at most
$14e$. For a nonperiod, its true endpoints lie in $Q_{3B},Q_{5B}$,
so \eqref{eq:patch-margin} supplies a witness of defect at least $\eta$.
All witnesses centered in $Q_B$ are selected. Restricting targets cannot
lower the true infimum. Its measured defect is at least $\eta-14e$.
Reserve another $6e$ for numerical enclosures and use threshold $\eta/2$,
with $20e<\eta/4$. This accepts exactly the true periods. Strict cutoff
collars ensure that every generator in Lemma~\ref{lem:recognition}
is present. The same tests on pairs of sources give their primitive
orbit relation: a period translating the first center to the second
must translate the first body to the second body, since Steiner points
lie in the interiors and different components are disjoint. All
primitive orbits are represented, so their number is recovered.

Accepted true vectors have a prior bound $U_0$ on their Euclidean
lengths. A nonzero determinant of two such periods is an integer
multiple of $\covol\Pi\ge V_*$. Determinant perturbations are
$O(U_0e)$; hence independence is decided below this gap. Select an
independent true pair with matrix $D$. Its generated sublattice has
index
\[
       n=|\det D|/\covol\Pi\le Q:=\max(1,\lceil U_0^2/V_*\rceil).
\]
Every coordinate of $D^{-1}w_j$ has reduced denominator dividing $n$
and a prior-bounded numerator. The inverse-matrix identity bounds its
estimation error by $C_*e$. Distinct rationals of denominators at most
$Q$ are separated by at least $Q^{-2}$. Finite search locks the exact
coordinates when $C_*e<1/(3Q^2)$.

Let $q$ be their common denominator, which divides $n$. A column
Hermite basis $H_0$ of the integer span of $qI_2$ and $qD^{-1}w_j$
gives
\begin{equation}\label{eq:period-basis}
                         \Pi=(DH_0/q)\Z^2.
\end{equation}
The integer group contains $q\Z^2$, so its Hermite entries have a
prior bound. Replacing $D$ by its estimate gives a matched true basis
with error $O(e)$. We have recovered exact integer relations, not
exact real coordinates. The primitive free area is
\[
             A_* =\covol\Pi-\sum_{j=1}^{r_*}\area(C_j).
\]
Uniform perimeter bounds and the parallel-body inclusions give area
error $O(e)$ for each matched component, proving the last assertion.
Support suprema and center integrals used in the tests admit certified
finite angular meshes, because the bounded bodies have a uniform
support Lipschitz bound. Their numerical errors are charged above.
\end{proof}

\begin{proof}[Proof of Theorem~\ref{thm:adaptive}]
Choose a sufficiently large fixed target square so that every component
needed in Lemma~\ref{lem:locking} and every coarse bracket is protected.
Separation bounds their number. Choose the fixed observation and launch
apertures as in Lemma~\ref{lem:stopped}. Apply
Proposition~\ref{prop:patch} with an integer angular mesh satisfying
$h\asymp\nu^{1/(s-2)}$ and a sufficiently small constant factor.
Its reconstructed bodies have matched $C^2$ support error $O(\nu)$
and Hausdorff error $e=O(h^s)$. For small $\nu$, all thresholds in
Lemma~\ref{lem:locking} hold. It gives exact primitive discrete data
and period-basis error $O(h^s)$. Repeating one recovered body per
primitive orbit under this basis gives a physical periodic presentation.
Indeed the inverse basis remains bounded, so only a bounded number of
integer translates can enter a bounded fundamental patch. The positive
separation and curvature margins then persist. The reconstructed
bodies have area error $O(h^s)$ and the covolume has the same order;
hence the primitive free-area error is also $O(h^s)$, in particular
$O(\nu)$.

The query count of Proposition~\ref{prop:patch} and its fixed local
stencil size give \eqref{eq:centers}. Equation~\eqref{eq:patch-cost}
gives \eqref{eq:attempts}, since $s$ is fixed. The calibration bias
allocation in Proposition~\ref{prop:query} gives \eqref{eq:precision}.
Every preparation has already been charged there. All steps are finite,
and the output is a finite smooth representation with certified
numerical tolerances. This proves the theorem.
\end{proof}

Without a known $\eta$, successive refinements give the corresponding
pointwise eventual decisions for each fixed table, but not a uniformly
valid observable stopping certificate. The exact functional inverse
combined with Lemma~\ref{lem:recognition} needs no supplied positive
margin: it tests zeros exactly. This distinction, as well as the
bounded periodic presentation itself, is retained. No inference of
crystallinity from an arbitrary nonperiodic configuration is made.

\section{A physical packing and a binary-transcript bound}
\label{sec:lower}
The lower bound counts the table-dependent binary outputs, not the
precision of the controller's choices. It therefore applies to more
powerful binary command designs as well as the reciprocal protocol.
Its geometric construction has fixed periods and a fixed positive
patch margin, so period uncertainty is not the cause of the bound.

\begin{lemma}[A uniform physical packing]\label{lem:packing}
For $s=6+\beta$ there is a fixed class $\mathcal P\subset\mathcal B_\eta$
and constants $c_0,c_1>0$ such that, for all sufficiently small $h$,
$\mathcal P$ contains $2^m$ tables, where $m\ge c_0/h$, whose variable
component support functions have pairwise $C^2$ distance at least
$c_1h^{s-2}$. All have the same known primitive lattice and two
primitive component orbits. Their correspondences and pose can be
supplied without reducing this separation.
\end{lemma}
\begin{proof}
Fix the lattice $20\Z^2$. Put one variable component near the unit
disk at the origin, and a fixed disk of radius two at $(6,0)$, and
repeat both by this lattice. Choose a nonzero smooth function
$\psi$ supported in $(-1/3,1/3)$ with $\psi''(0)\ne0$.
Take angular centers $\theta_j$ in a fixed interval with mutual
spacing at least $h$, with $c/h\le m\le C/h$. For
$\omega\in\{0,1\}^m$, define the uncentered support of the variable
body by
\begin{equation}\label{eq:packing}
 p_\omega(\theta)=1+a h^s
                   \sum_{j=1}^m\omega_j\psi((\theta-\theta_j)/h),
\end{equation}
periodically extended, where $a>0$ is a fixed sufficiently small
number. At most one summand is nonzero at each point. Its derivative
of order $k\le6$ is bounded by $Ca h^{s-k}$. The sixth derivative
has uniformly bounded $\beta$-H\"older seminorm. Within a support
this follows by rescaling; between supports it follows from their
separation and the vanishing at their edges. The same estimate follows
by comparison to the intervening zero point for two arbitrarily close
points on different sides of an edge. Thus the $C^{6,\beta}$ bound
is uniform over all $h,m,\omega$.

For small $a$, $p_\omega+p_\omega''$ lies in a fixed positive interval.
The bodies are smooth and strictly convex, with all separation,
diameter, lattice and curvature bounds independent of $h$. Subtracting
their Steiner first harmonic also preserves the uniform support prior.
The two component species cannot be translated into each other:
their diameters lie in disjoint fixed intervals. Every period must
therefore send a variable component to another variable component,
forcing that period to lie in $20\Z^2$. Conversely every such vector
is a period. This proves primitiveness and the stated orbit count.

Center differences of the same species are true periods. For any
cross-species difference, the translated source species cannot match
a target of the other species with small support error, because their
widths differ by a fixed amount. A target of the same species has
center discrepancy bounded below by a fixed positive number: the
cross-species offset $(6,0)$, perturbed by the variable Steiner point,
stays a positive distance from $20\Z^2$. There are representatives
of both species in $Q_B$. These two alternatives give a uniform
positive defect for every nonperiod candidate in
\eqref{eq:patch-margin}. Hence one fixed $\eta>0$ works throughout
the construction. Taking the union of these tables over all small
$h$ defines a single class $\mathcal P$ with fixed numerical bounds.

If $\omega\ne\omega'$, choose an index where they differ and evaluate
the second derivative at its center. All other summands vanish there,
so the absolute difference is $a|\psi''(0)|h^{s-2}$. This proves the
separation in the laboratory support norm. The variable component,
lattice and fixed component may be identified in advance; none of
these disclosures identifies its bump vector. No analytic smoothness
bound is used or claimed for this packing.
\end{proof}

\begin{lemma}[Binary decision trees with randomized commands]
\label{lem:bits}
Let $M$ parameters have pairwise disjoint acceptable output sets.
A procedure with at most $N$ binary table-dependent outputs, arbitrary
adaptive inputs and parameter-independent randomization has average
success probability at most $2^N/M$ under the uniform distribution on
these parameters. The binary response probabilities may be arbitrary
functions of the parameter and chosen input.
\end{lemma}
\begin{proof}
Represent the controller randomization by an independent seed. Each
binary response can be generated by comparing another independent
uniform random variable with its conditional response probability.
This constructs the joint adaptive experiment for every parameter on
one common seed space. Fix all these independent seeds. Commands are
then functions only of the preceding binary transcript, and the output
is a function of a binary tree with at most $2^N$ leaves, after padding
short transcripts if necessary. An output at a leaf is acceptable for
at most one parameter, by disjointness. Thus at most $2^N$ parameters
can succeed for these fixed seeds. Average over the uniform parameter
and then integrate over the seeds. The bound follows. The numerical
values of selected commands give no additional leaves, since they are
already determined by the seed and transcript.
\end{proof}

\begin{proof}[Proof of Theorem~\ref{thm:lower}]
In Lemma~\ref{lem:packing} take
$h\asymp\nu^{1/(s-2)}$ with its constant large enough that the support
separation exceeds $2\nu$. The $\nu$-accuracy output sets are disjoint.
Lemma~\ref{lem:bits}, with $M=2^m$, gives average success at most
$2^{N-m}$. Uniform success at least $3/4$ therefore requires
$N\ge m+\log_2(3/4)\ge c\nu^{-1/(s-2)}$ for small $\nu$.
This argument allows more input control than the upper-bound procedure
and assumes perfect calibration. It remains valid for the latter
procedure with all-attempt bookkeeping. A protocol that reveals a
free-start test or a continuous impact location has a different output
alphabet and is not covered by the assertion.
\end{proof}

The loss in this theorem is the supplied laboratory-frame loss used by
the active controller. It is already enough to establish a matching
power for Theorem~\ref{thm:adaptive}, in the same loss. We do not silently
replace it by an unlabelled quotient loss in the lower bound, nor extend
a deterministic cap $N$ to an expected stopping-time budget without
another argument. No lower bound for the logarithms or for confidence
dependence is asserted.

For fixed confidence, on a nondegenerate bounded class containing
$\mathcal P$, the two results give
\begin{equation}\label{eq:matched-power}
 c\nu^{-1/(4+\beta)}\le N^*(\nu)
          \le C\nu^{-1/(4+\beta)}\log^2(C/\nu),
\end{equation}
where $N^*$ is the least uniform deterministic cap on attempted bits,
and constants in the accuracy definition are absorbed by rescaling
$\nu$. The matching power comes from one-dimensional boundary
smoothness measured two derivatives below the prior. The physical
work needed for the upper bound is that these boundary queries can
be implemented from the stipulated collision bits without additional
geometric outputs or uncharged free preparations.

\section{Information categories and relation to earlier inverses}
\label{sec:comparison}
Active boundary estimation is an essential comparison for the finite
theorem, not merely an analogy with passive billiard spectra.
Castro and Nowak~\cite{CN} study adaptive binary classification under
boundary-fragment smoothness and noise assumptions, with matching
orders for their excess-risk criterion. Their input selects a feature
point whose response is a class label. This directly available label
is not our collision bit: a solid start and a free miss have the same
output here. Neither their observation model nor their loss identifies
our finite periodic table. We use an independent physical packing and
binary-tree proof rather than importing a classification lower bound.

Locatelli, Carpentier and Kpotufe~\cite{LCK} give an adaptive strategy
for smooth decision boundaries without prior knowledge of smoothness
and noise parameters. Their membership-query setting and local
line-search/interpolation construction make the closest algorithmic
comparison. In contrast, our smoothness is known, our error is a $C^2$
geometry norm, and the boundary label must first be computed from
reciprocal pooled collision means. We claim neither a new general
bisection principle nor an advance on adaptation to unknown smoothness.
The independent reduction in Proposition~\ref{prop:query} is what
permits this established active-estimation strategy in the physical
collision experiment.

The earlier scalar inverse reconstructed a uniformly accurate occupation
on a full two-dimensional grid. Its constructive upper bound was
$C\nu^{-3}\log(C/(\nu\delta))$ attempts on the same bounded
$C^{6,\beta}$ class, using $\sigma\asymp\nu^{3/2}$ and compensated
support smoothing. That theorem remains in Supplement R. The present
procedure uses $O(h^{-1})$ angular lines, a fixed-depth dependency
stencil for each queried point, and the full $C^{6,\beta}$ regularity.
It has a different adaptive command design, not a sharper analysis of
the old uniform design. We have not proved a matching lower bound
for that old design or for uniform spatial preparations.

The resource improvement is accompanied by the explicit localization
and calibration cost in \eqref{eq:precision}. Exact data remain a pair
of functionals on controlled spatial inputs. The lower bound concerns
only their finite binary realizations; it does not assert a low-information
passive invariant. The angle/time/position tolerances shrink with
localization and are not learned from the collision outcomes. Apparatus
motion, precision in representing input locations, and arithmetic
cost are distinct resources.

For completeness the exact v31 manuscript is supplied unchanged as
Supplement R. It retains the arbitrary centered-law and general
nontrivial bounded-law inverses, the different-zero-set comparison,
the full fixed-aperture proof, rational interval enclosures, the
uniform-grid construction, and all earlier moment, marked-law and
count-fiber results in their original nested volumes. The more general
randomized law is an exact functional theorem; the finite local query
proved here uses the four-atom compass law. No continuous transition
law is discretized without an error budget. Rational intervals remain
conditional on forcing and survival bounds, not certificates of the
apparatus or physical prior.

The finite global conclusion also retains the known positive patch
margin and the bounded periodic presentation. Exact rational period
relations are distinct from approximate Euclidean generators. Removing
the known margin gives only pointwise eventual recognition, as described
in Section~\ref{sec:periods}. The result does not settle rigidity from
uniformly prepared count germs, exact analytic count germs, marked
length spectra or passive trajectories. These boundaries delimit the
information theorem without changing its topic or weakening any retained
mathematical assertion.

\subsection*{Supplementary material and assisted analysis}
Supplement R is the complete reviewed A2 v31 manuscript, supplied with
all its mathematical inputs, nested preceding volumes and original
provenance unchanged. It is not an independently published article.
The present proofs do not require a missing supplementary argument.
AI-assisted analysis contributed to the local-query reduction, adaptive
reconstruction and information bound. The named author is responsible
for checking the arguments and references before submission. Finite
arithmetic diagnostics and compilation are not continuum proof
certificates, apparatus validation or independent referee acceptance.
\begin{thebibliography}{99}
\raggedright
\bibitem{CN}
R. M. Castro and R. D. Nowak, \emph{Minimax bounds for active learning},
IEEE Trans. Inform. Theory \textbf{54} (2008), 2339--2353.
\href{https://doi.org/10.1109/TIT.2008.920189}{doi:10.1109/TIT.2008.920189}.
\bibitem{LCK}
A. Locatelli, A. Carpentier, and S. Kpotufe,
\emph{An adaptive strategy for active learning with smooth decision boundary},
Proceedings of Algorithmic Learning Theory, Proceedings of Machine
Learning Research \textbf{83} (2018), 547--571.
\bibitem{LawlerLimic}
G. F. Lawler and V. Limic, \emph{Random Walk: A Modern Introduction},
Cambridge Studies in Advanced Mathematics \textbf{123}, Cambridge
University Press, 2010.
\bibitem{Retained}
Q. Qi, \emph{Scalar collision laws and recognition of periodic dispersing
billiards}, A2 v31, October 4, 2026. Complete reviewed manuscript included
as Supplement R with its preceding nested submission volumes; not cited
as an independently published article.
\end{thebibliography}

\end{document}
